\setcounter{section}{4}

  \zwischenueberschrift{Soal latihan}

\inputaufgabe
{}
{

Misalkan

\mavergleichskette
{\vergleichskette
{ h(x_1 , \ldots , x_n)
}
{ = }{ \sum_{i = 1}^n a_ix_i
}
{  }{
}
{    }{
}
{   }{
}
}
{}{}{}
adalah suatu
\definitionsverweis {bentuk linear}{}{}
$\neq 0$ dan

\mavergleichskette
{\vergleichskette
{ Y
}
{ = }{ h^{-1}(c)
}
{  }{
}
{    }{
}
{   }{
}
}
{}{}{}
adalah hipermuka yang bersesuaian. Buktikan bahwa
\definitionsverweis {pemetaan Weingarten}{}{}
merupakan pemetaan nol.

}
{} {}

\inputaufgabe
{}
{

Misalkan

\mavergleichskettedisp
{\vergleichskette
{ Y
}
{ =} { \left\{  \left(  x   , \,   y  , \,  z                 \right)   \mid   x^2+y^2 = z^2 , \,  \left(  x   , \,   y  , \,  z                 \right)  \neq  \left(  0   , \,   0  , \,   0                 \right)       \right\}
}
{ \subseteq} { W
}
{ =} { \R^3 \setminus \{  \left(  0   , \,   0  , \,   0                 \right) \}
}
{ } {
}
}
{}{}{}
adalah kerucut standar, dan misalkan

\mavergleichskette
{\vergleichskette
{ P
}
{ = }{ \left(  x   , \,   y  , \,  z                 \right)
}
{ \in }{ Y
}
{    }{
}
{   }{
}
}
{}{}{.}
Misalkan

\mavergleichskettedisp
{\vergleichskette
{ \gamma(t)
}
{ =} { \left(  x\cos t-y\sin t   , \,   x\sin t+y\cos t  , \,  z                 \right)
}
{ } {
}
{ } {
}
{ } {
}
}
{}{}{.}

Catatan edisi: kurva pada sumber tidak selalu melalui P dan tidak selalu tetap berada pada kerucut. Rumus di atas menggantinya dengan rotasi standar pada bidang xy.
Tentukan

\mathdisp {\operatorname{lim}_{   t  \rightarrow  0   } \,  { \frac{  N( \gamma(t) ) -  N( \gamma(0) )  }{ t } }} { , }

dengan $N$ medan normal satuan yang diperoleh dari gradien $x^2+y^2-z^2$.

}
{} {}

\inputaufgabe
{}
{

Misalkan

\mavergleichskettedisp
{\vergleichskette
{ Y
}
{ =} { \left\{  \left(  x   , \,   y  , \,  z                 \right)   \mid   x^2+y^2 =1         \right\}
}
{ \subseteq} { \R^3
}
{ } {
}
{ } {
}
}
{}{}{}
adalah silinder standar, dan misalkan

\mavergleichskette
{\vergleichskette
{ P
}
{ = }{ \left(  1   , \,   0  , \,  z                 \right)
}
{ \in }{ Y
}
{    }{
}
{   }{
}
}
{}{}{.}
Misalkan

\mavergleichskettedisp
{\vergleichskette
{ \gamma(t)
}
{ =} { \left(  \cos  t   , \,   \sin  t  , \,  z                 \right)
}
{ } {
}
{ } {
}
{ } {
}
}
{}{}{.}
Tentukan

\mathdisp {\operatorname{lim}_{   t  \rightarrow  0   } \,  { \frac{  N( \gamma(t) ) -  N( \gamma(0) )  }{ t } }} { , }

dengan $N$ medan normal satuan yang diperoleh dari gradien $x^2+y^2-1$.

}
{} {}

\inputaufgabe
{}
{

Misalkan

\mavergleichskettedisp
{\vergleichskette
{ Y
}
{ =} { \left\{  \left(  x   , \,   y  , \,  z                 \right)   \mid   x^2+y^2 =1         \right\}
}
{ \subseteq} { \R^3
}
{ } {
}
{ } {
}
}
{}{}{}
dan misalkan

\mavergleichskette
{\vergleichskette
{ P
}
{ \in }{ Y
}
{  }{
}
{    }{
}
{   }{}
}
{}{}{}
adalah suatu titik. Buktikan bahwa
\definitionsverweis {pemetaan Weingarten}{}{}
$L_P$ tidak bijektif, tetapi juga bukan pemetaan nol.

}
{} {}

\inputaufgabe
{}
{

Buktikan bahwa suatu
\definitionsverweis {permukaan terdiferensialkan}{}{}

\mavergleichskette
{\vergleichskette
{ Y
}
{ \subseteq }{ \R^3
}
{  }{
}
{    }{
}
{   }{
}
}
{}{}{}
dapat memuat suatu garis

\mavergleichskette
{\vergleichskette
{ G
}
{ \subseteq }{ Y
}
{  }{
}
{    }{
}
{   }{
}
}
{}{}{}
yang vektor arahnya, jika dipandang sebagai vektor di
\definitionsverweis {ruang singgung}{}{}
$T_PY$ untuk suatu titik

\mavergleichskette
{\vergleichskette
{ P
}
{ \in }{ G
}
{  }{
}
{    }{
}
{   }{
}
}
{}{}{,}
tidak harus termasuk dalam
\definitionsverweis {kernel}{}{}
dari
\definitionsverweis {pemetaan Weingarten}{}{}
$L_P$.

}
{} {}

\inputaufgabe
{}
{

Misalkan

\mavergleichskette
{\vergleichskette
{ W
}
{ \subseteq }{ \R^n
}
{  }{
}
{    }{
}
{   }{
}
}
{}{}{}
\definitionsverweis {terbuka}{}{,}
\maabb {h} { W } { \R
} {}
adalah suatu
\definitionsverweis {fungsi yang terdiferensialkan secara kontinu sebanyak dua kali}{}{}
dan

\mavergleichskette
{\vergleichskette
{ Y
}
{ = }{ h^{-1}(c)
}
{  }{
}
{    }{
}
{   }{
}
}
{}{}{}
adalah
\definitionsverweis {serat}{}{}
di atas

\mavergleichskette
{\vergleichskette
{ c
}
{ \in }{ \R
}
{  }{
}
{    }{
}
{   }{
}
}
{}{}{,}
dengan $h$
\definitionsverweis {reguler}{}{}
di setiap titik pada $Y$.
Catatan edisi: sumber hanya mengasumsikan h kelas C1. Di sini hipotesis diperkuat menjadi h kelas C2 agar medan normal gradien dan pemetaan Weingarten yang ditanyakan terdefinisi sebagaimana pada kuliah.
Kita juga memandang $h$ sebagai pemetaan $\tilde{h}$ pada
\mathl{W \times \R^m}{}{,}
dengan $m$ variabel terakhir tidak muncul secara eksplisit dalam $\tilde{h}$. Misalkan

\mavergleichskettedisp
{\vergleichskette
{ Z
}
{ =} { \tilde{h}^{-1}(c)
}
{ =} { Y \times \R^m
}
{ \subseteq} { W \times \R^{m}
}
{ } {
}
}
{}{}{.}
Misalkan

\mavergleichskette
{\vergleichskette
{ Q
}
{ \in }{ Z
}
{  }{
}
{    }{
}
{   }{
}
}
{}{}{}
adalah suatu titik yang terletak di atas

\mavergleichskette
{\vergleichskette
{ P
}
{ \in }{ Y
}
{  }{
}
{    }{
}
{   }{
}
}
{}{}{.}
Bagaimana hubungan antara
\definitionsverweis {pemetaan-pemetaan Weingarten}{}{}
\mathkor {} {L_Q} {dan} {L_P} {?}

}
{} {}

\inputaufgabegibtloesung
{}
{

Misalkan
\mathkor {} {r} {dan} {R} {}
adalah bilangan real positif dengan

\mavergleichskette
{\vergleichskette
{ 0
}
{ < }{ r
}
{ < }{ R
}
{    }{
}
{   }{
}
}
{}{}{,}
dan kita tinjau torus
\zusatzklammer {tertanam} {} {}

\mavergleichskettedisp
{\vergleichskette
{ T
}
{ =} { \left\{ (x,y,z) \in \R^3  \mid   \left(  \sqrt{x^2+y^2} -R  \right)^2 +z^2  = r^2         \right\}
}
{ } {
}
{ } {
}
{ } {
}
}
{}{}{.}
\aufzaehlungdrei{Tentukan suatu
\definitionsverweis {medan normal satuan}{}{}
$N$ pada $T$.
}{Pada titik
\mathl{\left( R-r  , \,   0  , \,   0                 \right)}{} dan untuk vektor

\mavergleichskette
{\vergleichskette
{ v
}
{ = }{ \begin{pmatrix}  0  \\ 1 \\  0   \end{pmatrix}
}
{  }{
}
{    }{
}
{   }{
}
}
{}{}{}
tentukan
\definitionsverweis {turunan berarah}{}{}

\mathl{D_vN}{.}
}{Pada titik
\mathl{\left( R-r  , \,   0  , \,   0                 \right)}{} dan untuk vektor

\mavergleichskette
{\vergleichskette
{ w
}
{ = }{ \begin{pmatrix}  0  \\ 0\\  1   \end{pmatrix}
}
{  }{
}
{    }{
}
{   }{
}
}
{}{}{}
tentukan
\definitionsverweis {turunan berarah}{}{}

\mathl{D_w N}{.}
}

}
{} {}

\inputaufgabe
{}
{

Untuk grafik $Y$ dari fungsi

\mavergleichskettedisp
{\vergleichskette
{ f(x,y)
}
{ =} { xy
}
{ } {
}
{ } {
}
{ } {
}
}
{}{}{}
tentukan
\definitionsverweis {pemetaan Weingarten}{}{}
$L_P$ pada setiap titik

\mavergleichskettedisp
{\vergleichskette
{ P
}
{ =} { (x,y,f(x,y) )
}
{ } {
}
{ } {
}
{ } {
}
}
{}{}{.}
Tentukan nilai eigen, vektor eigen, dan suatu matriks diagonal untuk $L_P$.

}
{} {}

\inputaufgabe
{}
{

Misalkan

\mavergleichskette
{\vergleichskette
{ W
}
{ \subseteq }{ \R^n
}
{  }{
}
{    }{
}
{   }{
}
}
{}{}{}
\definitionsverweis {terbuka}{}{,}
\maabb {h} { W } { \R
} {}
adalah suatu
\definitionsverweis {fungsi yang terdiferensialkan secara kontinu sebanyak dua kali}{}{}
dan

\mavergleichskette
{\vergleichskette
{ Y
}
{ = }{ h^{-1}(c)
}
{  }{
}
{    }{
}
{   }{
}
}
{}{}{}
adalah
\definitionsverweis {serat}{}{}
di atas

\mavergleichskette
{\vergleichskette
{ c
}
{ \in }{ \R
}
{  }{
}
{    }{
}
{   }{
}
}
{}{}{,}
dengan $h$
\definitionsverweis {reguler}{}{}
di setiap titik pada $Y$.
Catatan edisi: sumber hanya mengasumsikan h kelas C1. Di sini hipotesis diperkuat menjadi h kelas C2 agar pemetaan Weingarten yang dibandingkan terdefinisi sebagaimana pada kuliah.
Misalkan
\maabbdisp {L} {\R^n } { \R^n
} {}
adalah suatu
\definitionsverweis {isometri linear}{}{,}

\mavergleichskette
{\vergleichskette
{ W'
}
{ = }{ L^{-1}(W)
}
{  }{
}
{    }{
}
{   }{
}
}
{}{}{}
dan

\mavergleichskettedisp
{\vergleichskette
{ Z
}
{ =} { L^{-1}(Y)
}
{ =} { (h \circ L)^{-1}(c)
}
{ } {
}
{ } {
}
}
{}{}{.}
Buktikan bahwa
\definitionsverweis {pemetaan Weingarten}{}{}
$L_P$ untuk

\mavergleichskette
{\vergleichskette
{ P
}
{ \in }{ Y
}
{  }{
}
{    }{
}
{   }{
}
}
{}{}{}
bersesuaian dengan pemetaan Weingarten untuk

\mavergleichskette
{\vergleichskette
{ Q
}
{ = }{ L^{-1}(P)
}
{ \in }{ Z
}
{    }{
}
{   }{
}
}
{}{}{,}
jika melalui $L$ ruang singgung $T_PY$ dan $T_QZ$ dihubungkan satu sama lain.

}
{} {}

\inputaufgabegibtloesung
{}
{

Untuk permukaan $Y$ yang diberikan oleh

\mavergleichskettedisp
{\vergleichskette
{ x^2+y^2+2z^2
}
{ =} { 4
}
{ } {
}
{ } {
}
{ } {
}
}
{}{}{}
dan titik

\mavergleichskettedisp
{\vergleichskette
{ P
}
{ =} { \left(  1  , \,   1  , \,   1                 \right)
}
{ \in} { Y
}
{ } {
}
{ } {
}
}
{}{}{}
tentukan suatu matriks diagonal untuk
\definitionsverweis {pemetaan Weingarten}{}{}
$L_P$.

}
{} {}

\inputaufgabe
{}
{

Misalkan

\mavergleichskette
{\vergleichskette
{ W
}
{ \subseteq }{ \R^n
}
{  }{
}
{    }{
}
{   }{
}
}
{}{}{}
\definitionsverweis {terbuka}{}{,}
\maabb {h} { W } { \R
} {}
adalah suatu
\definitionsverweis {fungsi yang terdiferensialkan secara kontinu sebanyak dua kali}{}{}
dan

\mavergleichskette
{\vergleichskette
{ Y
}
{ = }{ h^{-1}(c)
}
{  }{
}
{    }{
}
{   }{
}
}
{}{}{}
adalah
\definitionsverweis {serat}{}{}
di atas

\mavergleichskette
{\vergleichskette
{ c
}
{ \in }{ \R
}
{  }{
}
{    }{
}
{   }{
}
}
{}{}{,}
dengan $h$
\definitionsverweis {reguler}{}{}
di setiap titik pada $Y$.
Bagaimana
\definitionsverweis {pemetaan Weingarten}{}{}
dan
\definitionsverweis {ruang-ruang eigennya}{}{}
berubah jika kita beralih dari suatu
\definitionsverweis {orientasi}{}{}
pada $Y$ ke orientasi yang berlawanan?

}
{} {}

  \zwischenueberschrift{Soal untuk dikumpulkan}

\inputaufgabe
{4}
{

Dalam
Contoh 4.5,
tentukan matriks representasi
\definitionsverweis {pemetaan Weingarten}{}{}
untuk titik-titik berbentuk
\mathl{\begin{pmatrix}  0  \\ y \\ z   \end{pmatrix}}{} terhadap basis
\mathkor {} {\begin{pmatrix}  y  \\ 0 \\  0   \end{pmatrix}} {dan} {\begin{pmatrix}  0  \\ z \\  y   \end{pmatrix}} {}
dari ruang singgung. Tentukan nilai-nilai eigen dan suatu basis yang terdiri atas vektor-vektor eigen.

}
{} {}

\inputaufgabe
{5}
{

Untuk permukaan $Y$ yang diberikan oleh

\mavergleichskettedisp
{\vergleichskette
{ x^2+2y^2+5z^2
}
{ =} { 11
}
{ } {
}
{ } {
}
{ } {
}
}
{}{}{}
dan titik

\mavergleichskettedisp
{\vergleichskette
{ P
}
{ =} { \left(  2  , \,   1  , \,   1                 \right)
}
{ \in} { Y
}
{ } {
}
{ } {
}
}
{}{}{}
tentukan suatu matriks diagonal untuk
\definitionsverweis {pemetaan Weingarten}{}{}
$L_P$.

}
{} {}

\inputaufgabe
{4}
{

Untuk hipermuka yang diberikan oleh

\mavergleichskettedisp
{\vergleichskette
{ x^2+2y^2+3z^2+4w^2
}
{ =} { 10
}
{ } {
}
{ } {
}
{ } {
}
}
{}{}{}

\mavergleichskette
{\vergleichskette
{ Y
}
{ \subseteq }{ \R^4
}
{  }{
}
{    }{
}
{   }{
}
}
{}{}{}
dan titik

\mavergleichskettedisp
{\vergleichskette
{ P
}
{ =} { \left(  1  , \,   1  , \,   1 , \,   1                \right)
}
{ \in} { Y
}
{ } {
}
{ } {
}
}
{}{}{}
tentukan suatu matriks untuk
\definitionsverweis {pemetaan Weingarten}{}{}
$L_P$ terhadap suatu
\definitionsverweis {basis}{}{}
yang sesuai dari
\definitionsverweis {ruang singgung}{}{}
$T_PY$.

}
{} {}

\inputaufgabe
{6}
{

Untuk grafik $Y$ dari fungsi

\mavergleichskettedisp
{\vergleichskette
{ f(x,y)
}
{ =} { x^3y^5
}
{ } {
}
{ } {
}
{ } {
}
}
{}{}{}
tentukan
\definitionsverweis {pemetaan Weingarten}{}{}
$L_P$ pada titik-titik

\mavergleichskette
{\vergleichskette
{ P_1
}
{ = }{ (0,0,0)
}
{  }{
}
{    }{
}
{   }{
}
}
{}{}{,}

\mavergleichskette
{\vergleichskette
{ P_2
}
{ = }{ (0,1,0)
}
{  }{
}
{    }{
}
{   }{
}
}
{}{}{,}

\mavergleichskette
{\vergleichskette
{ P_3
}
{ = }{ (2,1,8)
}
{  }{
}
{    }{
}
{   }{
}
}
{}{}{.}
Untuk masing-masing titik, tentukan nilai eigen, vektor eigen, dan suatu matriks diagonal untuk $L_P$.

}
{} {}
